\section*{Chapter 14 --- Bus analysis}
\addcontentsline{toc}{section}{Chapter 14 --- Bus analysis}

\solution{ex:14:1}{Read your own traffic}
The exercise asks for measurements; the questions are answered here.
\ans{a} They should match exactly what you wrote to \code{TX_ID}, \code{TX_DLC} and the data
registers. If they do not, the fault is in the packing, not in the transfer --- the analyser
decodes what is actually on the wires.
\ans{b} The order is what the analyser adds. A swap of \code{TX_DATA_LO} and \code{TX_DATA_HI}
gives a frame with the right identifier and the right length but the wrong byte order, and this is
where it first becomes visible: every test on the laptop may have been green if the test made the
same mistake as the code.
\ans{c} Among others SOF, RTR, IDE, r0, the CRC field, the ACK slot and EOF. The register map
exposes only the identifier, the DLC and the data --- the controller handles the rest, and the
analyser is the only place where you see them.

\solution{ex:14:2}{The other way round}
The exercise asks for measurements; the questions are answered here.
\ans{a} Yes, if RX valid was set and you read the registers. If it does not: check that the
identifier is one your node receives, and that the bit rate is the same in the analyser.
\ans{b} Bit 1 should go high when the frame is received and low only when you write \code{0x1} to
\code{RX_ACK}. If it stays set, the acknowledgement has not gone through --- check that bit 0
really is set in the written word.
\ans{c} The second frame overwrites the first if you have not had time to acknowledge it. The
controller holds exactly one received frame, and the newest wins; that is documented and
deliberate. With a \code{receive()} taking around 270~µs and frames taking 130~µs, you do not have
time.

\solution{ex:14:3}{The acknowledgement}
\ans{a} Yes. The analyser is a real CAN controller and pulls the ACK slot low on every frame it
receives with a correct CRC. In the trace it shows as the frame being registered without an
acknowledgement error.
\ans{b} The error flag remains \emph{untouched}. This design never reads the ACK slot back, so it
cannot detect that nobody acknowledged --- the transmission completes exactly as if everything had
gone well.
\ans{c} A complete CAN controller would have detected the acknowledgement error, incremented its
transmit error counter and \textbf{retransmitted the frame}, over and over until somebody answered
or the counter reached bus off. A lone node on a real CAN bus therefore transmits the same frame
for ever. This controller does not, because it has neither ACK read-back nor automatic
retransmission --- two of the documented simplifications.

\solution{ex:14:4}{The bottleneck}
\ans{a} It depends on how often your application calls \code{send()}. With one frame per iteration
in a tight loop the load is low by CAN standards --- a few per cent --- since the SPI link slows
you down long before the bus fills up.
\ans{b} A \code{send()} is six SPI transactions of around 45~µs each, that is, about
\textbf{270~µs}. A maximal CAN frame is around 130 bits, that is, about \textbf{130~µs} at
1~Mbit/s.
\ans{c} \textbf{The SPI link.} It takes roughly twice as long as the frame it sends, so the CAN
bus stands waiting for you. That matches the calculation in \secref{sec:timing} --- and this is
where you can see with your own eyes that it did, which is why the chapter asks you to measure
instead of trusting the table.
